% Terminal apparatus of the concise study: the scope note and the References
% for the sources this companion actually uses. The revision timestamp and the
% rights colophon follow from the shared generation record.

\section{Appendix: Scope and Qualifications}

{\footnotesize
\textbf{Edition, formulary and occurrence.} This concise study treats the Mass
\latin{Dominica vigesima post Pentecosten}, II classis, of the \work{Missale
Romanum}, editio typica, Typis Polyglottis Vaticanis 1962, pp.~412--413,
nos.~1689--1698, in the universal 1962 calendar and no particular calendar. Its
occurrence on 11~October 2026, its precedence over the feast of the Maternity
of the Blessed Virgin Mary, the commemoration of that feast in low and
conventual Masses and its omission in a sung Mass that is not conventual, the
colour, the Credo and the Trinity Preface rest on the Missal's own calendar and
general rubrics (RG 16, 108, 111, 127~b; RGMR 434, 437~a, 475~a, 494~b). The
commemorated orations, which the full study prints, are not studied here.

\textbf{Relation to the full study.} Every claim here is drawn from the full
study of this same formulary, the Full PDF subtitled \textit{The word that
heals from afar, the song of the exiles, and the two peoples: a study of the
proper in three readings}, which prints the appointed Latin with its
public-domain English, treats each element in its own setting, and develops the
three readings at length, each with its own four senses; nothing
evidence-dependent is added here. The four rows on the first page orient a
reader and do not replace those readings' own senses. The map and the
dossier's explanatory rows condense the study's. Each Father and later writer
is cited only for what he said of his own passage, and each liturgical
commentator only for what he says of the elements of this formulary he names.

\textbf{Text, English and rights.} Every Latin form of the ten elements was read
for the study on the page images of the CMAA facsimile of the typical edition
and collated with the Pustet \work{Missale Romanum} of Ratisbon, 1862, in which
each element's wording already stands; the collation is in
\texttt{propers/verified.md}. The Introit's antiphon is a cento of words from
Dan 3:29--43 that no witness read carries in the Missal's form, and the Latin
source of the Gradual's \latin{illis escam}, the Alleluia's \latin{psallam
tibi, gloria mea} and the Communion's added \latin{Domine} was not found in the
Vulgate, the Greek Psalter or the Latin Fathers read; no Roman or Old Latin
psalter and no chant manuscript was consulted. Psalms
are cited in the Missal's Vulgate numbering, with the Hebrew numbering added in
the map and the dossier. The scriptural English is the Douay--Rheims in
Challoner's revision at the canonical verses, so that it does not render the
Missal's Latin where a chant departs from its verse; the English of the
Introit's antiphon and of the three orations is the anonymous translation
printed by Cummiskey in Philadelphia in 1861. Neither is an approved liturgical
translation of the 1962 Missal. The English of St Augustine, St John
Chrysostom and St Cyril of Alexandria is that of the Nicene and Post-Nicene
Fathers and the Library of the Fathers; English beside the Latin of St Gregory,
St Jerome, St Hilary and St Augustine is a gloss of that Latin. The Latin of the
formulary rests on its public-domain antecedent, the Pustet Missal (17 U.S.C.
103(b)); the Douay--Rheims, the 1861 English, Migne's texts, the
nineteenth-century translations of the Fathers, O'Sullivan's English of
Bellarmine, the 1883 English of the continuation of \work{The Liturgical Year}
and the 1927 English of Schuster's \work{Sacramentary} are in the public domain
in the United States. St Gregory's Latin is quoted from the Latin Wikisource
transcription of the \work{Homiliarum in Evangelia}, licensed under Creative
Commons Attribution-ShareAlike 3.0, with attribution to its contributors at
\url{https://la.wikisource.org/wiki/Homiliarum_in_Evangelia}. Cassiodorus in
the \work{Corpus Christianorum} edition and the \work{New American Bible}
introductions are described or summarised, not reproduced.

\textbf{Reception and its limits.} The witnesses named here are those the
study's sweep read at the loci below. St Augustine's expositions of Pss~118,
136 and 144, St Hilary's tractates and St Jerome's commentary on Ephesians were
read where their Latin is quoted on page images of Migne's \work{Patrologia
Latina}; St Gregory's homily in the Wikisource transcription of Migne's text.
St Thomas Aquinas, Cassiodorus, St John Chrysostom's Greek exposition of
Ps~136, Sicard and Durandus were read in text layers or a restricted scan and
are described, not quoted. Among the witnesses not reached are Origen's
commentary on St John, St Ambrose on Ps~118, Theodoret and St Athanasius on the
psalms, Theophylact on Ephesians, the
Breviary's homily for this Sunday and St Alphonsus Liguori's Sunday sermons.

\textbf{Chronology and records.} The Date cells on the second page come from
the Scripture chronology corpus through the record generated for this
formulary, with one declared comparison under its critical profile for the
Epistle; the explanatory rows say which source each date comes from and how it
is qualified. The search, its loci, the disagreements and the negative results
are in \texttt{research/scope.md}, the reasoning behind the three readings in
\texttt{research/interpretations.md}, and this companion's production record in
\texttt{research/production-review.md}.
\par}

\section*{References}
\runninghead{References}

{\scriptsize
\subsection*{Liturgical books and Bibles}
\begin{itemize}\setlength{\itemsep}{0.15em}\setlength{\parskip}{0pt}
\item \work{Missale Romanum ex decreto Sacrosancti Concilii Tridentini restitutum, Summorum Pontificum cura recognitum}, editio typica (Typis Polyglottis Vaticanis, 1962), CMAA facsimile, \url{https://media.churchmusicassociation.org/pdf/missale62.pdf}: \work{Rubricae generales} 16, 108, 111, 127; \work{Rubricae generales Missalis} 434, 437, 475, 494; \work{Ordo Missae} no.~1033; pp.~412--413 (nos.~1689--1698).
\item \work{Missale Romanum} (Ratisbon: Pustet, 1862), pp.~352--354, \latin{Dominica XX. post Pentecosten}.
\item \work{The Roman Missal, Translated into the English Language for the Use of the Laity} (Philadelphia: Eugene Cummiskey, 1861), pp.~448--450 (XX Sunday after Pentecost, Introit, Collect, Secret, Postcommunion).
\item The Holy Bible, Douay--Rheims version, revised by Bishop Richard Challoner: Deut 32:15; 2 Mac 15:14; Ps 107:2, 4; 118:1, 19, 49--50; 136:1, 4, 9; 144:15--17; Isa 44:4; Jer 25:11--12; Dan 3:1, 20, 24--45; Mic 7:1; Mt 8:7; Jn 4:41--54; Rom 11:25--26; 1 Cor 14:22; Eph 1:1; 4:1; 5:8, 15--21; 6:20. \work{Biblia Sacra Vulgatae editionis} (Clementine Vulgate): Ps 107:2; 118:1, 49--50; 136:1, 4; 144:15--16; Dan 3:29--43; Jn 4:46--53; Eph 5:15--21.
\end{itemize}

\subsection*{Fathers, Doctors and saints}
\begin{itemize}\setlength{\itemsep}{0.15em}\setlength{\parskip}{0pt}
\item St Augustine, \work{Tractates on the Gospel of John} 16, 3--5, 7, Nicene and Post-Nicene Fathers, 1st ser., vol.~7. \work{Enarrationes in Psalmos}, 1st ser., vol.~8: on Ps~56:8; 118, sermons 1 and 15 (on vv.~1, 49--50); 136, on vv.~1--2, 9; 144, on vv.~15--17. Latin in Migne, PL 37 (Paris, 1845), cols.~1502 (118, sermo 1, 1), 1761 and 1774 (136), 1882 (144).
\item St Gregory the Great, \work{Homiliae in Evangelia} 4, 3 and 28, 1--3, Latin Wikisource transcription of Migne's text, CC BY-SA 3.0; Migne, PL 76, cols.~1210--1213.
\item St Jerome, \work{Commentarii in Danielem}, on Dan 3:26--49, in Migne, PL 25 (Latin Wikisource transcription), cols.~639--642; \work{Commentarii in Epistolam ad Ephesios} III, on Eph 5:16--20, in Migne, PL 26 (Paris, 1845), cols.~527--530.
\item St Hilary of Poitiers, \work{Tractatus super Psalmos}, in Ps.~118, \latin{Aleph} 1--3 and \latin{Zain}; in Ps.~136, in Migne, PL 9 (Paris, 1844), cols.~503--505, 548--550, 777--784.
\item St John Chrysostom, \work{Homilies on St John} 35, 2--3, Nicene and Post-Nicene Fathers, 1st ser., vol.~14 (New York, 1889), pp.~123--125; \work{Homilies on Ephesians} 19, 1st ser., vol.~13; \work{Expositio in Psalmos}, in Ps.~136, Migne, PG 55, cols.~405--407 (described, not quoted).
\item St Cyril of Alexandria, \work{Commentary on the Gospel according to S.~John} II.5, trans. P.~E. Pusey, Library of the Fathers (Oxford, 1874), pp.~232--236.
\item St Thomas Aquinas, \work{Super Evangelium S.~Ioannis lectura} c.~4, lect.~7 (Venice, 1745), pp.~521--524; \work{In Epistolam ad Ephesios} c.~5, lect.~6 (Liège: H.~Dessain, 1857), pp.~344--345 (described, not quoted).
\item St Robert Bellarmine, \work{A Commentary on the Book of Psalms}, trans. J.~O'Sullivan: on Pss~56:8; 118:1, 49; 136.
\item Cassiodorus, \work{Expositio Psalmorum}, on Pss~107:2, 4; 118, preface and v.~1; 136; 144:15, ed. M.~Adriaen, CCSL 98 (Turnhout, 1958) (described, not quoted).
\end{itemize}

\subsection*{Liturgical commentators and chronology sources}
\begin{itemize}\setlength{\itemsep}{0.15em}\setlength{\parskip}{0pt}
\item Bl.\ Ildefonso Schuster, \work{The Sacramentary (Liber Sacramentorum)}, vol.~3 (London: Burns Oates \& Washbourne, 1927), pp.~175--177.
\item The continuation of Prosper Guéranger, \work{The Liturgical Year} (Dom Lucien Fromage), \work{The Time after Pentecost}, vol.~2 (series vol.~11), trans. Laurence Shepherd (Dublin: James Duffy and Sons, 1883), pp.~439, 447, 451--452.
\item Sicard of Cremona, \work{Mitrale} VIII.20, in Migne, PL 213; William Durandus, \work{Rationale divinorum officiorum} VI.137 (Lyon, 1612) (text layers; described, not quoted).
\item \work{The Catholic Encyclopedia} (New York, 1907--1912): ``Captivities of the Israelites'' (vol.~3, 1908); ``Book of Daniel'' and ``David'' (vol.~4, 1908); ``Epistle to the Ephesians'' (vol.~5, 1909); ``Jesus Christ'', ``The Gospel of Saint John'', ``Third and Fourth Books of Kings'' and ``Chronology of the Kings'' (vol.~8, 1910); ``St.\ Paul'' (vol.~11, 1911); ``The New Testament'' and ``Temple of Jerusalem'' (vol.~14, 1912).
\item G.~L. Haydock, \work{The Douay--Rheims Bible with Haydock's Commentary}, notes on Pss~70 and 136.
\item United States Conference of Catholic Bishops, \work{New American Bible, Revised Edition}, introductions to the Psalms and to Ephesians, on their dating; summarised, not reproduced.
\end{itemize}
\par}
